% 09_computational_proof.tex
% Ch.9: Computational proof + anomaly catalog

\chapter{Computational Proof and Anomaly Catalog}
\label{ch:computational_proof}

\begin{objectivebox}
\begin{itemize}
    \item Defend the computational universality of the AVE physics engine ($39$ orders of magnitude) without free parameters.
    \item Establish the strict numerical guard hierarchy ($10^{-12}$ to $10^{-30}$) ensuring topological mathematics never cross-contaminate with machine-precision floating-point artefacts.
    \item Prove that the $r_n \propto n^2$ orbital scaling law does not rely on Bohr or Schrödinger postulates, but emerges natively from the geometry of continuous acoustic impedance cavities acting on discrete matter.
\end{itemize}
\end{objectivebox}

\section{Computational Proof of Scale Invariance}

The full computational proof is implemented in the physics engine module
\texttt{src/ave/}\allowbreak\texttt{axioms/}\allowbreak\texttt{scale\_}\allowbreak\texttt{invariant.py}. Every domain module in the
framework (\texttt{saturation.py}, \texttt{fdtd\_3d.py}, and 13 other domain modules) imports its impedance
operations from this single canonical source. The automated test suite
verifies that the same
function produces identical numerical results whether called with
particle-scale, molecular-scale, laboratory-scale, geophysical-scale,
or galactic-scale inputs.

This is not an analogy. It is a \textbf{structural identity}: the same
operator, the same code, the same mathematics, from $10^{-13}$~m to
$10^{26}$~m.

\section{Verification Summary}

\begin{center}
\small
\begin{tabularx}{\textwidth}{@{}l l l r X@{}}
\toprule
\textbf{Domain} & \textbf{Prediction} & \textbf{Regime} & \textbf{Agreement} & \textbf{Params} \\
\midrule
Galactic rotation (NGC 3198)  & $v = 159$ km/s         & I/IV  & 5\%    & 0 \\
Multi-galaxy RAR              & McGaugh curve           & I--IV & exact  & 0 \\
'Oumuamua acceleration        & $4.54 \times 10^{-6}$ m/s$^2$ & I & 91\%   & 0 \\
Kirkwood gaps (5/5)           & $< 0.3\%$ error         & I     & $<$0.3\% & 0 \\
Earth magnetopause            & 9.1 R$_E$               & I     & 8.7\%  & 0 \\
Jupiter magnetopause          & 55.6 R$_J$              & I     & 11.8\% & 0 \\
Neutrino MSW $P_{ee}$ (4 channels) & 0.31--0.55        & I     & $<$10\% & 0 \\
Superconductor $B_c(T)$ (5 materials) & $\sqrt{1-(T/T_c)^2}$ & II/III & exact & 0 \\
London depth ($\lambda_L$)    & 37--150 nm             & II/III & exact  & 0 \\
Seismic $\Gamma_{\text{Moho}}$ & 0.29                  & I     & matches PREM & 0 \\
GW lossless propagation       & $V_{GW}/V_{snap} = 10^{-28}$ & I & exact & 0 \\
Topological Pair Production   & $H_{net}=0 \to e^+ + e^-$    & IV    & exact & 0 \\
Protein folding (CLN025)      & RMSD = 2.59 \AA        & I     & sub-3 \AA & 0 \\
\bottomrule
\end{tabularx}
\end{center}

\begin{figure}[h]
    \centering
    \includegraphics[width=1.0\textwidth]{cross_scale_verification.png}
    \caption{The single Axiom~4 saturation kernel $S(x, x_{yield}) = \sqrt{1 - (x/x_{yield})^2}$ governs physics from the lattice pitch ($10^{-13}$~m) to the Hubble radius ($10^{26}$~m) --- 39 orders of magnitude with the same operator, the same code, zero adjustable parameters.}
    \label{fig:cross_scale_verification}
\end{figure}

\section{Anomaly Catalog: Proposed Tests}
\label{sec:anomaly_catalog}

The following anomalies and phenomena are identified as targets for
future AVE verification:

\subsection{Tier 2: Precision Tests}
\begin{enumerate}
    \item \textbf{Flyby anomaly:} Unexplained $\Delta v$ of $\sim$3--13 mm/s
    during Earth flybys. Predict: impedance gradient at magnetopause
    produces differential phase shift.
    \item \textbf{Muon $g-2$:} $\Delta a_\mu = 2.49 \times 10^{-9}$.
    AVE predicts a vacuum self-energy correction from lattice topology.
    \item \textbf{Hubble tension:} $H_0 = 67$ (CMB) vs 73 (Cepheid).
    AVE's $H_\infty$ is a single number; scale-dependent impedance
    may resolve the tension.
    \item \textbf{Proton radius puzzle:} Already solved --- $d_p =
    4\hbar/m_p c \approx 0.841$ fm matches muonic hydrogen.
\end{enumerate}

\subsection{Tier 3: Deep Predictions}
\begin{enumerate}
    \item \textbf{Fast radio bursts:} Dispersion measure via
    $Z_{plasma}(r)$ integration along line of sight.
    \item \textbf{Cosmic ray knee/ankle:} Impedance-driven GZK limit
    from lattice saturation at $V/V_{snap} \sim 1$.
    \item \textbf{Pioneer anomaly:} Solar wind drag vs thermal
    radiation pressure on the RTG.
    \item \textbf{Saturn ring system:} Ring gaps as impedance waveguide
    modes (cf. Kirkwood gaps).
    \item \textbf{Enceladus plume:} Saturn's Io analog --- tidal heating
    $\to$ cryovolcanism $\to$ plasma source.
\end{enumerate}

\subsection{Tier 4: Hardware Verification}
\begin{enumerate}
    \item \textbf{Chiral Antennas:} Wire-stitched impedance antenna measuring
    vacuum $Z_0$ directly.
    \item \textbf{High-Voltage Dielectric Rectification:} Asymmetrical solid-state capacitor in mineral oil bath for $\varepsilon_r$ sweep.
    \item \textbf{Casimir measurement:} Impedance between parallel
    plates as function of separation.
\end{enumerate}

\section{Numerical Precision and Dimensional Coordination}
\label{sec:precision_policy}

The physics engine enforces a rigorous numerical precision policy to ensure that computational artefacts never masquerade as physical predictions.

\subsection{Floating-Point Arithmetic}
All engine modules use IEEE~754 double precision (\texttt{float64}), providing approximately 15 significant decimal digits ($\varepsilon_{mach} \approx 2.2 \times 10^{-16}$).  No single-precision (\texttt{float32}) arrays are used anywhere in the framework.  This is critical because the AVE derivation chain spans 39 orders of magnitude: a constant computed at the lattice pitch ($\ell_{node} \sim 10^{-13}$~m) participates in predictions at the Hubble radius ($R_H \sim 10^{26}$~m), requiring at least 12 significant digits to remain unambiguous.

\subsection{Guard Constant Hierarchy}
Three standardised guard constants prevent numerical singularities throughout the engine.  All modules import these from \texttt{constants.py} rather than defining ad-hoc magic numbers:

\begin{center}
\small
\renewcommand{\arraystretch}{1.3}
\begin{tabular}{llp{7.2cm}}
\toprule
\textbf{Constant} & \textbf{Value} & \textbf{Usage} \\
\midrule
\texttt{EPS\_NUMERICAL} & $10^{-12}$ & Impedance ratios, reflection coefficients, normalisation denominators.  Ensures $Z/(Z + \varepsilon) \approx 1$ within float64 precision. \\
\texttt{EPS\_CLIP} & $10^{-15}$ & Saturation operator clip ceiling: $\sqrt{1 - x^2}$ with $x = 1 - \varepsilon$.  Prevents $\sqrt{<0}$ while keeping $S(A_{yield}) \approx 0$ to float64 resolution. \\
\texttt{EPS\_DIVZERO} & $10^{-30}$ & Hard division-by-zero floor for denominators that can reach exactly zero (DC impedance, $\tanh$ at origin).  Sub-float64 so it never affects physical results. \\
\bottomrule
\end{tabular}
\end{center}

These constants are \emph{dimensionless} ratios applied to already-normalised quantities.  They carry no units and encode no physics.

\subsection{Dimensional Analysis Chain}
Every derived constant in the engine is expressed algebraically from the three calibration inputs ($m_e$, $\alpha$, $G$) plus exact SI definitions ($c_0$, $\mu_0$, $\varepsilon_0$, $\hbar$, $e$).  The dimensional algebra is as follows:

\begin{resultbox}{Dimensional Traceability}
\begin{alignat*}{3}
\ell_{node} &= \hbar / (m_e c_0) \quad &&[\text{m}] \quad &&\text{(Input 1: spatial cutoff)} \\
\alpha &= e^2 / (4\pi\varepsilon_0 \hbar c_0) \quad &&[\text{---}] \quad &&\text{(Input 2: dielectric bound)} \\
G &= \hbar c_0 / (7 \xi_M m_e^2) \quad &&[\text{m}^3/\text{kg}\cdot\text{s}^2] \quad &&\text{(Input 3: Machian boundary)}
\end{alignat*}
\end{resultbox}

Every subsequent constant (Bohr radius, Rydberg energy, string tension, snap voltage, yield field, proton mass, etc.) is an explicit algebraic combination of these three with no hidden fitting.

\subsection{Precision Budget}

\begin{center}
\small
\renewcommand{\arraystretch}{1.2}
\begin{tabular}{lll}
\toprule
\textbf{Constant Class} & \textbf{Source} & \textbf{Relative Uncertainty} \\
\midrule
SI exact ($c_0$, $\mu_0$, $e$, $\hbar$, $k_B$) & 2019 SI redefinition & 0 (exact by definition) \\
$m_e$ (CODATA 2018) & Penning trap & $3.0 \times 10^{-10}$ \\
$\alpha$ (CODATA 2018) & Electron $g$-factor & $1.5 \times 10^{-10}$ \\
$G$ (CODATA 2018) & Torsion balance & $2.2 \times 10^{-5}$ \\
$m_p$ (CODATA 2018) & Penning trap & $3.1 \times 10^{-10}$ \\
\bottomrule
\end{tabular}
\end{center}

The gravitational constant $G$ is the least precise input at $2.2 \times 10^{-5}$ (22 ppm).  This uncertainty propagates only into cosmological and gravitational predictions (MOND $a_0$, Machian hierarchy $\xi_M$).  All electromagnetic, atomic, and molecular predictions are limited only by $m_e$ and $\alpha$ uncertainties below 1~ppb, which is far below the ODE solver tolerance and effectively exact.

\section{Avoidance of Methodological Contamination}
\label{sec:methodological_contamination}

A fundamental principle of the AVE engine is complete independence from external theoretical frameworks. Theoretical contamination occurs when a formula derived from an incompatible physical postulate (e.g., probability amplitudes, hidden geometry, or non-AVE empirical fits) is unknowingly imported as a ``convenient approximation.''

A prominent example discovered during framework development involved atomic ionisation energies. Early iterations of the solver attempted to correct standard shielding models, culminating in formulas of the form $IE = Z_{eff}^2 \, Ry / n^2$. While directionally correct, this explicitly imports the Schrödinger solution for a hydrogenic atom (or the ad-hoc Bohr model). It is not an AVE equation.

The framework requires all energy states to emerge directly and exclusively from the 5-step universal regime-boundary eigenvalue method. The Rydberg energy itself is not a postulate or a given parameter; it is an emergent consequence of the electron cavity saturation boundary ($r_{sat} = a_0 = \ell_{node}/\alpha$), yielding exact reproduction of the $13.6057$ eV constant through pure topological scaling:
\begin{equation}
    Ry = \left[ (\hbar c(1+\nu)) / a_0 \right] \times \frac{\alpha/(1+\nu)}{2} = \frac{\alpha^2 m_e c^2}{2}
\end{equation}

\subsection{Topological Orbital Radii vs. Quantum Mechanics}
Does the framework implicitly rely on QM by assuming discrete orbital shells scale as $r_n = n^2 a_0 / Z$? No. In AVE, discrete orbitals are macroscopic standing wave cavities — explicitly identical to the physics governing Saturn's rings and Kirkwood gaps in the asteroid belt. 

In the solar system, mean-motion orbital gaps occur where periods are commensurate: $(T/T_{moon})^2 = (q/p)^2$. This periodic boundary condition creates a Fabry-Pérot impedance cavity. At the atomic scale, the ``moon'' is the electron itself. Its orbital circumference must fit an integer number of its own spatial standing waves (the topological phase wind):
\begin{equation}
    2\pi r = n \lambda
\end{equation}
Since the local central impedance requires the phase velocity to scale as $v \propto 1/\sqrt{r}$, and impedance matching requires $\lambda \propto 1/v$, the wavelength scales as $\lambda \propto \sqrt{r}$. Substituting this into the resonance condition yields:
\begin{equation}
    2\pi r \propto n \sqrt{r} \implies \sqrt{r} \propto n \implies r_n \propto n^2
\end{equation}

Therefore, $r_n = n^2 a_0 / Z$ is not a QM assumption; it is the unique topological standing-wave condition for a scale-free central impedance profile. The Schrödinger equation merely discovers this same geometric constraint heuristically through operators.

When extending to multi-electron atoms, the engine must never revert to $Z_{eff}$ arithmetic; it must strictly apply the coupled resonator normal mode solver using universal pairwise energy terms to find the resulting eigenvalues.

\subsection{The Electron Saturation Radius}

\begin{examplebox}[Verifying the Linear Saturation Regime for Orbitals]
\textbf{Problem:} The topological saturation radius of the electron defines its structural limit as $d_{sat} = \ell_{node}$. Given that the macroscopic acoustic separation between atomic electrons is characterized by the Bohr radius ($a_0 = \ell_{node}/\alpha$), evaluate the pairwise voltage strain and classify its elastodynamic regime.

\textbf{Solution:} The pairwise voltage strain ($V/V_{snap}$) is the ratio of the saturation radius to the interaction distance ($A = d_{sat} / r_{ij}$). Substituting the structural equivalents:
\[ A = \frac{d_{sat}}{r_{ij}} \approx \frac{\ell_{node}}{\ell_{node}/\alpha} = \alpha \]
Evaluating the magnitude:
\[ \alpha \approx 0.007297 \]
Because $0.007 \ll \sqrt{2\alpha}$ (the $\approx 0.121$ boundary between Linear Regime I and Saturation Regime II), electron-electron repulsion operates deep within the perfectly linear, Hookean domain of the LC network. Non-linear saturated impedance corrections are entirely unnecessary; standard linear Coulomb electrostatics mapped onto topological geometries apply exactly.
\end{examplebox}

\subsection{Multi-Electron Repulsion: Topological Subshell Impedance Cascade}

In standard quantum mechanics, peer-electron Coulomb repulsion involves time-averaging spatial probability densities via $N$-body continuous integration. As discussed in Section~\ref{sec:methodological_contamination}, introducing such integrals into AVE is a methodological contamination that artificially recovers the macroscopic shielding collapse.

The framework requires a strictly scale-invariant, 0-parameter resolution. By treating atomic orbitals as nested macroscopic cavities rather than overlapping probability clouds, peer interactions are dictated by the \textbf{Subshell Impedance Cascade} (see Chapter~\ref{ch:quantum_orbitals} for the complete derivation). 

Instead of continuous integration, the discrete lattice geometry enforces:
\begin{enumerate}
    \item \textbf{Cross-Shell Gauss Screening (Axiom 2):} Outer solitons perceive exact stepwise integer screening $Z_{\rm net}(r) = Z - N_{\rm inner}$.
    \item \textbf{Same-Shell Topologic Node Sorting (Axiom 4):} Electrons within the same principal shell ($n$) couple exclusively as discrete geometric LC resonators subject to the $K = 2G$ topological limit and the $8\pi\alpha$ spatial packing fraction. 
\end{enumerate}

This pure discrete formalism flawlessly bounds multi-electron systems without resorting to empirical variables, $Z_{\rm eff}$ fitting, or continuous variational probability integrals. The ``multi-electron correlation problem'' of standard quantum mechanics is bypassed entirely by enforcing topological geometry.

\section*{Chapter Summary}
\begin{itemize}
    \item The AVE Physics Engine operates identically from the $10^{-13}\text{ m}$ lattice pitch to the $10^{26}\text{ m}$ Hubble boundary, definitively proving the structural scale-invariance of the underlying continuum mathematics.
    \item Theoretical contamination via empirically fitted scaling laws (e.g., probability amplitudes or parameterized fractional charges) is rigorously barred from the engine.
    \item All constants inside the model are purely algebraically traceable back to the triad of fundamental calibration boundaries ($m_e$, $\alpha$, $G$).
    \item Macroscopic multi-body repulsion derives analytically from the Topologic Subshell Impedance cascade (Phase D Architecture), entirely bypassing the "N-body correlation" integration crisis of QM.
\end{itemize}

\section*{Exercises}
\begin{enumerate}
    \item Why must the physics engine execute exclusively in \texttt{float64} precision, and why is an artificial floor (like \texttt{EPS\_DIVZERO} = $10^{-30}$) required when computing perfectly antipodal discrete topological interference metrics?
    \item In deriving the Hydrogen energy levels, AVE establishes that the quantization integer $n$ does not originate from abstract quantum boundary conditions, but instead from the macro-acoustic resonance rule ($2\pi r = n \lambda$). Why is a continuous macroscopic standing wave structurally impossible if $n = 1.5$?
\end{enumerate}
